\documentclass[tikz,border=10pt]{standalone}
\usepackage{tikz}
\usepackage[fontset=none]{ctex}
\setCJKmainfont{Noto Sans CJK SC}
\setCJKsansfont{Noto Sans CJK SC}
\usetikzlibrary{shapes, arrows.meta, positioning, fit, backgrounds, calc, shadows, decorations.pathmorphing}

\definecolor{drawBlueFill}{HTML}{DAE8FC}
\definecolor{drawBlueLine}{HTML}{6C8EBF}
\definecolor{drawGreenFill}{HTML}{D5E8D4}
\definecolor{drawGreenLine}{HTML}{82B366}
\definecolor{drawOrangeFill}{HTML}{FFE6CC}
\definecolor{drawOrangeLine}{HTML}{D79B00}
\definecolor{drawPurpleFill}{HTML}{E1D5E7}
\definecolor{drawPurpleLine}{HTML}{9673A6}
\definecolor{drawRedFill}{HTML}{F8CECC}
\definecolor{drawRedLine}{HTML}{B85450}
\definecolor{drawGreyFill}{HTML}{F5F5F5}
\definecolor{drawGreyLine}{HTML}{666666}

\pgfdeclarelayer{background}
\pgfsetlayers{background,main}

\begin{document}
\begin{tikzpicture}[
    every node/.style={font=\footnotesize},
    base_box/.style={rectangle, rounded corners=3pt, align=center,
        minimum height=1.0cm, minimum width=2.4cm,
        drop shadow={opacity=0.15}, thick},
    blue_node/.style={base_box, fill=drawBlueFill, draw=drawBlueLine},
    green_node/.style={base_box, fill=drawGreenFill, draw=drawGreenLine},
    orange_node/.style={base_box, fill=drawOrangeFill, draw=drawOrangeLine},
    purple_node/.style={base_box, fill=drawPurpleFill, draw=drawPurpleLine},
    red_node/.style={base_box, fill=drawRedFill, draw=drawRedLine},
    arrow/.style={-{Stealth[scale=1.2]}, thick, color=black!70},
    arrow_data/.style={-{Stealth[scale=1.2]}, very thick, color=drawOrangeLine},
    arrow_dash/.style={-{Stealth[scale=1.2]}, dashed, thick, color=black!70},
    phase_title/.style={font=\bfseries\small, text=black!80},
    icon_circle/.style={circle, minimum size=0.9cm, thick, drop shadow={opacity=0.1}, align=center},
]

% ===== 阶段标题 =====
\node[phase_title] (t1) at (0, 1.0) {数据准备};
\node[phase_title] (t2) at (4.2, 1.0) {特征工程};
\node[phase_title] (t3) at (8.4, 1.0) {模型训练};
\node[phase_title] (t4) at (12.6, 1.0) {评估部署};

% ===== 阶段一 =====
\node[icon_circle, fill=drawBlueFill, draw=drawBlueLine] (ico1) at (0, -0.2) {};
\draw[drawBlueLine, thick, fill=drawBlueFill!60]
    (-0.2, -0.05) -- (-0.2, -0.35) arc (180:360:0.2 and 0.08) -- (0.2, -0.05);
\draw[drawBlueLine, thick] (-0.2, -0.05) arc (180:360:0.2 and 0.08);
\draw[drawBlueLine, thick] (-0.2, -0.05) arc (180:0:0.2 and 0.08);

\node[blue_node, below=0.6cm of ico1] (collect) {多源数据采集};
\node[blue_node, below=0.5cm of collect] (clean) {清洗与标注};

% ===== 阶段二 =====
\node[icon_circle, fill=drawGreenFill, draw=drawGreenLine] (ico2) at (4.2, -0.2) {};
\draw[drawGreenLine, thick] (4.2, -0.2) circle (0.22);
\draw[drawGreenLine, thick] (4.2, -0.2) circle (0.10);
\foreach \a in {0,45,...,315} {
    \draw[drawGreenLine, thick] (4.2, -0.2) ++ (\a:0.22) -- ++(\a:0.08);
}

\node[green_node, below=0.6cm of ico2] (extract) {特征提取};
\node[green_node, below=0.5cm of extract] (select) {特征选择与降维};

% ===== 阶段三 =====
\node[icon_circle, fill=drawOrangeFill, draw=drawOrangeLine] (ico3) at (8.4, -0.2) {};
\draw[drawOrangeLine, thick, fill=drawOrangeFill!60, rounded corners=1pt]
    (8.15, -0.35) rectangle (8.65, -0.05);
\foreach \x in {8.22, 8.32, 8.42, 8.52} {
    \draw[drawOrangeLine, thick] (\x, -0.05) -- (\x, 0.0);
    \draw[drawOrangeLine, thick] (\x, -0.35) -- (\x, -0.4);
}

\node[orange_node, below=0.6cm of ico3] (train) {分布式训练};
\node[orange_node, below=0.5cm of train] (hyper) {超参数调优};

% ===== 阶段四 =====
\node[icon_circle, fill=drawPurpleFill, draw=drawPurpleLine] (ico4) at (12.6, -0.2) {};
\draw[drawPurpleLine, thick, fill=drawPurpleFill!60]
    (12.6, 0.05) -- (12.45, -0.25) -- (12.5, -0.25) -- (12.5, -0.4)
    -- (12.7, -0.4) -- (12.7, -0.25) -- (12.75, -0.25) -- cycle;

\node[purple_node, below=0.6cm of ico4] (eval) {模型评估};
\node[purple_node, below=0.5cm of eval] (deploy) {容器化部署};

% ===== 主流程箭头 =====
\draw[arrow_data] (clean.east) -- ++(0.4,0) |- (extract.west);
\draw[arrow_data] (select.east) -- ++(0.4,0) |- (train.west);
\draw[arrow_data] (hyper.east) -- ++(0.4,0) |- (eval.west);

% 阶段内部
\draw[arrow] (collect) -- (clean);
\draw[arrow] (extract) -- (select);
\draw[arrow] (train) -- (hyper);
\draw[arrow] (eval) -- (deploy);

% ===== 反馈回路（落在 extract 节点左侧边，不落在图标上）=====
\draw[arrow_dash, color=drawRedLine]
    (eval.north) -- ++(0, 0.5) -- ($(extract.north west)+(0, 0.5)$) -- (extract.north west);
\node[font=\scriptsize, text=drawRedLine, fill=white, inner sep=1pt]
    at ($(eval.north)!0.5!(extract.north west) + (0, 0.5)$) {效果不佳时回退};

% ===== 底部指标 =====
\node[red_node, minimum width=3.0cm, below=1.2cm of deploy] (metric) {准确率 / 延迟 / 吞吐量};
\draw[arrow] (deploy) -- (metric);

% ===== 阶段分区 =====
\begin{pgfonlayer}{background}
    \node[fill=drawBlueFill!10, draw=drawBlueLine!30, dashed, rounded corners=8pt,
          inner sep=12pt, fit=(ico1)(collect)(clean)] {};
    \node[fill=drawGreenFill!10, draw=drawGreenLine!30, dashed, rounded corners=8pt,
          inner sep=12pt, fit=(ico2)(extract)(select)] {};
    \node[fill=drawOrangeFill!10, draw=drawOrangeLine!30, dashed, rounded corners=8pt,
          inner sep=12pt, fit=(ico3)(train)(hyper)] {};
    \node[fill=drawPurpleFill!10, draw=drawPurpleLine!30, dashed, rounded corners=8pt,
          inner sep=12pt, fit=(ico4)(eval)(deploy)] {};
\end{pgfonlayer}

\end{tikzpicture}
\end{document}
